\documentclass[conference]{IEEEtran}

\usepackage[T1]{fontenc}
\usepackage[utf8]{inputenc}
\usepackage[spanish,es-tabla]{babel}
\usepackage{lmodern}
\usepackage{microtype}
\usepackage{csquotes}
\usepackage{amsmath,amssymb}
\usepackage{booktabs}
\usepackage{array}
\usepackage{tabularx}
\usepackage{url}
\usepackage[hidelinks]{hyperref}

% Gestion bibliografica mediante biblatex con backend biber (Estilo IEEE)
\usepackage[backend=biber,style=ieee,sorting=none]{biblatex}
\addbibresource{referencias.bib}

\setlength{\parindent}{0.65em}
\setlength{\parskip}{1pt}
\renewcommand{\arraystretch}{1.18}
\DefineBibliographyStrings{spanish}{bibliography={Referencias}}
\urlstyle{same}

\title{Taller de Identificaci\'on y Clasificaci\'on de Tipos de Requisitos:\\Caso Plataforma de Gesti\'on de Citas M\'edicas}

\author{
  \IEEEauthorblockN{Jos\'e Ignacio Gilces Carranza}
  \IEEEauthorblockA{Universidad T\'ecnica Estatal de Quevedo (UTEQ) --- Facultad de Ciencias de la Computaci\'on\\
  Ingenier\'ia de Software $\cdot$ Ingenier\'ia de Requisitos (ISR-401) $\cdot$ Actividad Formativa FORM-TL-01}
}

\begin{document}
\onecolumn
\maketitle

\begin{abstract}
En este informe se presenta el desarrollo formal de la actividad FORM-TL-01, enfocada en la clasificaci\'on tridimensional (tipo de requisito e ISO/IEC 25010:2023, \'ambito producto/proceso y frontera sistema/software) de los 20 enunciados del caso Policl\'inico Ribera del Quevedo. Se realiza la reescritura y cuantificaci\'on rigurosa de 5 requisitos de calidad vagos mediante sintaxis can\'onica ISO/IEC/IEEE 29148:2018, definiendo m\'etricas, unidades y umbrales verificables alineados con la demanda operativa local y la LOPDP del Ecuador.
\end{abstract}

\section{Contexto del Sistema y Stakeholders}
El Policl\'inico Ribera del Quevedo atiende a 12\,000 pacientes registrados con un promedio de 450 citas diarias en 9 especialidades. La saturaci\'on telef\'onica y el ausentismo del 22\,\% exigen una plataforma web y m\'ovil integrada a un quiosco de autoservicio y al sistema de facturaci\'on preexistente.

\begin{center}
\footnotesize
\begin{tabularx}{\textwidth}{@{}p{3.2cm}p{3.8cm}>{\raggedright\arraybackslash}X@{}}
\toprule
\textbf{Stakeholder} & \textbf{Rol en el Policl\'inico} & \textbf{Necesidad preliminar de elicitaci\'on}\\
\midrule
Pacientes y familiares & Usuarios finales del portal/m\'ovil & Reservar, reprogramar o cancelar citas sin filas y recibir recordatorios por SMS/correo.\\
Recepcionistas & Operadores de ventanilla & Registrar la llegada presencial, reducir la transcripci\'on manual y atender sin saturaci\'on.\\
M\'edicos especialistas & Usuarios del m\'odulo m\'edico & Consultar el flujo diario de pacientes en consultorio y marcar atenciones realizadas.\\
Director M\'edico & Gestor operativo y patrocinador & Reducir el ausentismo del 22\,\%, optimizar horarios y obtener reportes de inasistencias.\\
Administrador & Gestor de configuraci\'on & Gestionar cat\'alogo de m\'edicos, especialidades, feriados y supervisar el flujo de cobros.\\
\'Area de TI (2 personas) & Custodios de infraestructura & Mantener el sistema sobre Java 21 y PostgreSQL con alta disponibilidad.\\
Proveedor de Facturaci\'on & Entidad externa de integraci\'on & Proveer la interfaz de servicios web para autorizar e intercambiar cobros.\\
\bottomrule
\end{tabularx}
\end{center}

\section{Clasificaci\'on Tridimensional de Requisitos}
Se clasifican los 20 enunciados del sistema (E01 a E20) conforme a tres dimensiones: (1) Tipo de requisito e ISO/IEC 25010:2023 \cite{iso25010}, (2) \'Ambito (Producto vs. Proceso) y (3) Frontera (Sistema vs. Software) \cite{pohl2025,iso29148}.

\begin{center}
\scriptsize
\setlength{\tabcolsep}{3pt}
\begin{tabularx}{\textwidth}{@{}p{0.6cm}p{3.6cm}p{3.4cm}p{1.4cm}p{1.4cm}>{\raggedright\arraybackslash}X@{}}
\toprule
\textbf{ID} & \textbf{Enunciado del Caso} & \textbf{Tipo / ISO 25010} & \textbf{\'Ambito} & \textbf{Frontera} & \textbf{Justificaci\'on T\'ecnica}\\
\midrule
E01 & Reserva de cita por especialidad, m\'edico, fecha y hora. & Funcional & Producto & Software & Especifica una funci\'on directa de interacci\'on del paciente con el sistema.\\
E02 & Env\'io de recordatorio SMS 24 horas antes de la cita. & Funcional & Producto & Software & Comportamiento aut\'onomo del software ante un evento temporal programado.\\
E03 & La aplicaci\'on debe ser r\'apida. & Calidad (Eficiencia de desempe\~no) & Producto & Software & Atributo de calidad del software: requiere cuantificaci\'on de tiempos de respuesta.\\
E04 & Registro de llegada del paciente en ventanilla por recepcionista. & Funcional & Producto & Software & Funci\'on operativa asignada al rol recepcionista en la interfaz.\\
E05 & Desarrollo en Java 21 con base de datos PostgreSQL. & Restricci\'on (T\'ecnica / Stack) & Producto & Software & Limita las decisiones de arquitectura al stack dominado por TI.\\
E06 & Consulta de pacientes del d\'ia y marcado de atenci\'on por m\'edico. & Funcional & Producto & Software & Funci\'on de gesti\'on de flujo de atenci\'on en el consultorio m\'edico.\\
E07 & Los datos de salud de los pacientes deben estar seguros. & Calidad (Seguridad / Confidencialidad) & Producto & Software & Exige mecanismos de protecci\'on de datos sensibles en el software.\\
E08 & Funcionamiento en Android 10 o superior. & Restricci\'on (Entorno operativo) & Producto & Software & Condici\'on de compatibilidad de ejecuci\'on sobre el SO m\'ovil.\\
E09 & Cancelaci\'on con $<$ 4h se registra como inasistencia. & Restricci\'on (Regla de negocio) & Producto & Sistema & Regla institucional del policl\'inico que gobierna la l\'ogica del servicio.\\
E10 & F\'acil de usar para pacientes adultos mayores. & Calidad (Usabilidad / Operabilidad) & Producto & Software & Caracter\'istica de interacci\'on y accesibilidad enfocada en el usuario final.\\
E11 & Entrega del proyecto en un plazo m\'aximo de 6 meses. & Restricci\'on (Proyecto / Gesti\'on) & Proceso & Sistema & L\'imite temporal que condiciona el ciclo de vida del desarrollo.\\
E12 & Cierre de sesi\'on autom\'atico tras 10 min de inactividad. & Restricci\'on (Regla de seguridad) & Producto & Software & Mecanismo de control de sesi\'on para prevenir accesos no autorizados.\\
E13 & Lectura de c\'edula en pantalla e impresi\'on de ticket en quiosco. & Funcional & Producto & Sistema & Involucra hardware (pantalla/impresora), firmware y software del quiosco.\\
E14 & La plataforma debe estar siempre disponible. & Calidad (Fiabilidad / Disponibilidad) & Producto & Software & M\'etrica de continuidad operativa del servicio web/m\'ovil.\\
E15 & Generaci\'on de reporte mensual de inasistencias por especialidad. & Funcional & Producto & Software & Funci\'on de procesamiento anal\'itico de datos para la direcci\'on m\'edica.\\
E16 & Validaci\'on de incrementos con Director M\'edico cada 2 semanas. & Restricci\'on (Metodol\'ogica) & Proceso & Sistema & Define el marco de revisi\'on y aceptaci\'on en el proceso de desarrollo.\\
E17 & Cumplimiento de la Ley Org\'anica de Protecci\'on de Datos (LOPDP). & Restricci\'on (Legal / Regulatoria) & Producto & Sistema & Marco normativo nacional de obligatorio cumplimiento.\\
E18 & Integraci\'on sin problemas con el sistema de facturaci\'on actual. & Calidad (Compatibilidad / Interoperabilidad) & Producto & Software & Capacidad del software para intercambiar datos con el sistema externo.\\
E19 & Restablecimiento en $<$ 30 min sin p\'erdida tras fallo de servidor. & Calidad (Fiabilidad / Recuperabilidad) & Producto & Software & Especifica el RTO ($<$30 min) y RPO (0 p\'erdidas) ante contingencias.\\
E20 & Revisi\'on de c\'odigo por pares antes de integrar en rama principal. & Restricci\'on (Calidad de proceso) & Proceso & Software & Pr\'actica obligatoria de aseguramiento de calidad en el repositorio Git.\\
\bottomrule
\end{tabularx}
\end{center}

\section{Cuantificaci\'on de Requisitos de Calidad Vagos}
Los enunciados vagos E03, E07, E10, E14 y E18 se reescriben aplicando la estructura can\'onica de la norma ISO/IEC/IEEE 29148:2018 (``El sistema deber\'a [acci\'on] [objeto] cuando [condici\'on], con el fin de [prop\'osito]'') y las caracter\'isticas de calidad de ISO/IEC 25010:2023 \cite{iso29148,iso25010}.

\vspace{0.25cm}
\begin{center}
\scriptsize
\setlength{\tabcolsep}{3pt}
\begin{tabular}{@{}p{0.7cm}p{5.1cm}p{2.3cm}p{1.4cm}p{0.8cm}p{2.8cm}p{2.8cm}@{}}
\toprule
\textbf{ID} & \textbf{Requisito Cuantificado (ISO 29148)} & \textbf{M\'etrica} & \textbf{Umbral} & \textbf{Unidad} & \textbf{Condiciones de Medici\'on} & \textbf{M\'etodo / Herramienta}\\
\midrule
\textbf{E03-Q} & El sistema deber\'a responder las consultas de disponibilidad de citas en 2,5 segundos o menos en el percentil 95 cuando operen 300 usuarios concurrentes durante 30 minutos sobre 12\,000 registros, con el fin de evitar retrasos en horas pico. & Percentil 95 del tiempo de respuesta HTTP & $\le$ 2{,}5 & s & 300 usuarios virtuales; 30 min de carga; servidor en homologaci\'on. & Prueba de carga automatizada con Apache JMeter.\\
\addlinespace[3pt]
\textbf{E07-Q} & El sistema deber\'a aplicar cifrado AES-256 en reposo y denegar el 100\,\% de 300 solicitudes de acceso no autorizado a datos de salud cuando se ejecuten pruebas de autorizaci\'on, con el fin de proteger la confidencialidad seg\'un la LOPDP \cite{lopdp2021}. & Tasa de denegaci\'on de accesos indebidos & 100\% & \% & 300 peticiones de prueba con cuentas sin permisos requeridos; AES-256. & Inspecci\'on y pruebas de autorizaci\'on con Postman.\\
\addlinespace[3pt]
\textbf{E10-Q} & El sistema deber\'a permitir que al menos el 85\,\% de 20 adultos mayores ($>$ 60 a\~nos) completen la reserva de una cita en 180 segundos o menos sin asistencia, utilizando interfaz con fuente $>$ 16 pt y contraste WCAG AAA en Android 10, con el fin de garantizar la accesibilidad. & Tasa de finalizaci\'on de tarea sin asistencia & $\ge$ 85\%; $\le$ 180 & \% / s & 20 usuarios mayores de 60 a\~nos; tarea de agendamiento m\'ovil sin ayuda previa. & Prueba de usabilidad con observaci\'on directa y cron\'ometro.\\
\addlinespace[3pt]
\textbf{E14-Q} & El sistema deber\'a mantener una disponibilidad operativa de al menos el 99,2\,\% en la plataforma web y m\'ovil durante el horario de atenci\'on (lunes a s\'abado 07:00--19:00) evaluado mensualmente, con el fin de garantizar el acceso continuo a turnos. & Porcentaje de tiempo de servicio activo (Uptime) & $\ge$ 99{,}2\% & \% mensual & Monitoreo continuo mediante sondas sint\'eticas HTTP cada 60 s durante 30 d\'ias. & An\'alisis h\'ibrido de logs con Uptime Kuma.\\
\addlinespace[3pt]
\textbf{E18-Q} & El sistema deber\'a sincronizar correctamente el 100\,\% de 200 solicitudes de verificaci\'on de pago con el sistema de facturaci\'on actual mediante su interfaz web documentada con una latencia inferior a 1,5 s, con el fin de evitar cobros duplicados. & Tasa de sincronizaci\'on de transacciones y latencia & 100\%; latencia $<$ 1{,}5 & \% / s & 200 peticiones de cobro simuladas contra la API del proveedor en homologaci\'on. & Pruebas de integraci\'on de contrato con SoapUI / Postman.\\
\bottomrule
\end{tabular}
\end{center}
\vspace{0.25cm}

\noindent\textbf{Justificaci\'on de umbrales.} Los 300 usuarios concurrentes responden al pico m\'aximo previsto los lunes a las 07:00. El umbral de 2,5~s asegura fluidez en dispositivos Android de gama media. La tasa del 100\,\% en seguridad y compatibilidad es requerida por tratarse de datos sensibles bajo la LOPDP \cite{lopdp2021} e integraci\'on financiera. El 85\,\% de usabilidad en 180~s considera la curva de aprendizaje de adultos mayores (35\,\% de la poblaci\'on del policl\'inico). La disponibilidad del 99,2\,\% tolera un m\'aximo de 2,3 horas mensuales de indisponibilidad programada.

\section{Referencias}
\printbibliography[heading=none]

\end{document}